\begin{thebibliography}{99}
\raggedright
\addcontentsline{toc}{section}{References}

\bibitem{ProveItSnapshot}
V.~Reshetnikov and contributors, \emph{ProveIt: Analysis / Polylogarithms},
manuscript and research reports, repository snapshot of 10 October 2026,
commit \texttt{a2a4cf58c49c745058c40e4a6748d472a3420f18}.
\href{https://github.com/VladimirReshetnikov/ProveIt/tree/a2a4cf58c49c745058c40e4a6748d472a3420f18/Analysis/Polylogarithms/docs/manuscript}{Pinned manuscript snapshot on GitHub}.
The accompanying source manifest records the inspected chapter blobs.

\bibitem{PriorResearch}
ProveIt incoming research report,
\path{polylogarithms_research_2026-10-10.zip}, at the snapshot above.
In particular, the Gaussian section proves the residual level-three
and level-four polynomial descriptions used here.
\href{https://github.com/VladimirReshetnikov/ProveIt/tree/a2a4cf58c49c745058c40e4a6748d472a3420f18/Analysis/Polylogarithms/docs/incoming}{Incoming reports at the pinned snapshot}.

\bibitem{PriorRigidity}
ProveIt incoming research report,
\path{polylogarithms_rigidity_20261010.zip}, at the snapshot above.
In particular, the Gaussian section gives the Möbius-duality certificate
for $S_4$. The archive's Git blob is recorded in the incoming manifest.

\bibitem{Zhao2010}
J.~Zhao, \emph{Standard relations of multiple polylogarithm values at
roots of unity}, Documenta Mathematica \textbf{15} (2010), 1--34.
\url{https://arxiv.org/abs/0707.1459}.

\bibitem{Panzer2017}
E.~Panzer, \emph{The parity theorem for multiple polylogarithms},
Journal of Number Theory \textbf{172} (2017), 93--113.
\url{https://arxiv.org/abs/1512.04482}.

\bibitem{Umezawa2025}
R.~Umezawa, \emph{An explicit parity theorem for multiple polylogarithms},
arXiv:2508.02040 (2025).
\url{https://arxiv.org/abs/2508.02040}.

\bibitem{RadchenkoZagier}
D.~Radchenko and D.~Zagier,
\emph{Arithmetic properties of the Herglotz function},
Journal für die reine und angewandte Mathematik \textbf{797} (2023),
229--253.
\url{https://doi.org/10.1515/crelle-2023-0009}.
\url{https://arxiv.org/abs/2012.15805}.
Author manuscript:
\url{https://people.mpim-bonn.mpg.de/zagier/files/preprints/HerglotzFunction.pdf}.

\bibitem{ChoieKumar2023}
Y.~Choie and R.~Kumar,
\emph{Arithmetic properties of the Herglotz--Zagier--Novikov function},
Advances in Mathematics \textbf{433} (2023), Article 109315.
\url{https://doi.org/10.1016/j.aim.2023.109315}.
\url{https://arxiv.org/abs/2309.10634}.

\bibitem{DixitSathyanarayanaSharan2024}
A.~Dixit, S.~Sathyanarayana, and N.~Guru Sharan,
\emph{Mordell--Tornheim zeta functions and functional equations for
Herglotz--Zagier type functions}, arXiv:2405.07934 (2024).
\url{https://arxiv.org/abs/2405.07934}.

\bibitem{DLMFGamma}
NIST Digital Library of Mathematical Functions, Section 5.8,
\emph{Infinite products}, reciprocal-gamma Weierstrass product.
\url{https://dlmf.nist.gov/5.8}. Accessed 10 October 2026.

\bibitem{Lindemann1882}
F.~Lindemann, \emph{Ueber die Zahl $\pi$},
Mathematische Annalen \textbf{20} (1882), 213--225.
\url{https://doi.org/10.1007/BF01446522}.

\bibitem{Kuznetsov2019}
A.~Kuznetsov, \emph{Lagrange inversion theorem for Dirichlet series},
arXiv:1911.11133 (2019).
\url{https://arxiv.org/abs/1911.11133}.

\end{thebibliography}
